\section{Finite polynomial factorizations give original native projectivity}
\label{sec:native-projectivity-from-polynomial-factorization}

\begin{theorem}[Original projectivity over an affine coefficient spectrum]
\label{thm:native-projectivity-from-polynomial-factorization}
\lean{EnvelopeProperty.NativeProjectivityFromPolynomialFactorization.polynomialProjection_isProjective, EnvelopeProperty.NativeProjectivityFromPolynomialFactorization.of_isProjectiveFactorization, EnvelopeProperty.NativeProjectivityFromPolynomialFactorization.iff_isProjectiveFactorization}
\leanok
\uses{thm:symmetric-polynomial-basis-comparison, thm:polynomial-projective-coefficient-isomorphism}
\source{bj22-trivval:sec:tc}
For every commutative ring $R$, every original finite polynomial
projective projection is projective in the original native
finite-quasicoherent-projective-bundle sense. Consequently an original
morphism to $\operatorname{Spec}R$ is native projective if and only if
it admits a finite polynomial closed factorization, over that same
original ring. No alternate base or supplied comparison certificate
is used.
\end{theorem}

\begin{proof}
Let $Y=\operatorname{Spec}R$, and let $E$ be the actual finite free
sheaf with $n+1$ generators on $Y$. Its original global sections are
free of rank $n+1$ over $\Gamma(Y,\mathcal O_Y)$: evaluation preserves
the finite biproduct defining the free sheaf, and the evaluated unit
sheaf is exactly that original coefficient ring. Transport the
standard product basis along this actual biproduct isomorphism.
The preceding basis theorem identifies the entire original symmetric
projective bundle of these sections with the polynomial projective
space over $\Gamma(Y,\mathcal O_Y)$. The whole affine projective-bundle
comparison identifies this with the entire original $E$-bundle.
Finally apply the actual original isomorphism
$\Gamma(\operatorname{Spec}R,\mathcal O)\simeq R$ to the coefficients.
Its spectrum is exactly the canonical affine comparison map; both
whole coefficient diagrams therefore show that the constructed
isomorphism retains the original projection to the same
$\operatorname{Spec}R$. This proves native projectivity of the
polynomial projection using $E$ itself. Compose any original closed
immersion into that polynomial space with the constructed actual
bundle isomorphism. Closed immersions compose, proving native
projectivity of the original morphism. The reverse implication is
the already proved actual finite polynomial presentation of every
finite quasicoherent projective bundle over an affine scheme.
\end{proof}
